\chapter{Profiling y diagnóstico}
\label{chap:profiling}

\section{Objetivo del profiling}
\label{sec:profiling-objetivo}

Los capítulos anteriores establecen dos hechos. El modo~3 reduce de forma
marcada los tiempos frente al modo~2, y solo EP y FT obtienen aceleración
sostenida. El objetivo de este capítulo es relacionar esos resultados con
métricas de ejecución de bajo nivel, como la ocupación de CPU, el IPC y los
eventos de caché. El diagnóstico se apoya en dos herramientas. La primera es
\texttt{perf stat}~\cite{LinuxPerfStat}, que mide el uso de recursos a nivel de
proceso. La segunda es el nuevo sistema de muestreo de Python~3.15
(Sección~\ref{sec:profiling-py315}), que atribuye el tiempo de sincronización y
el balance de carga por función e hilo.

El profiling con \texttt{perf stat} mide el proceso completo y, por tanto, no
aísla únicamente la ventana temporizada por el benchmark NAS. En este capítulo se
usa una matriz común para todos los benchmarks, con Python~3.15t, clase~W, modo~3 y
1, 8 y 32~hilos. Así, la lectura de CPU, IPC y caché no mezcla efectos de versión
con diferencias de diseño de la campaña. Cada medición incluyó una ejecución de
calentamiento previa para reducir el efecto de la compilación Cython sobre los
contadores.

Las métricas de \texttt{perf stat} deben interpretarse con cautela en
benchmarks cortos. IS clase~W tarda centésimas de segundo. En ese caso, el
arranque del proceso y las bibliotecas importadas pueden pesar tanto como el
kernel. También se observan valores de CPUs utilizadas superiores a uno en
algunas ejecuciones con un hilo, precisamente porque \texttt{perf} mide el
proceso completo y no solo la región paralela.


\section{Control de entorno y afinidad}
\label{sec:profiling-afinidad}

Las campañas finales no fijan por defecto las variables de afinidad de OpenMP
(\texttt{OMP\_PROC\_BIND}, \texttt{OMP\_PLACES} y \texttt{GOMP\_CPU\_AFFINITY}),
y las eliminan de forma explícita en los scripts de ejecución salvo en
experimentos controlados. El motivo es un efecto que conviene precisar, porque
no nace de OMP4Py. La biblioteca no fija hilos a núcleos. Lee
\texttt{OMP\_PROC\_BIND} a una variable interna sin aplicarla y la cláusula
\texttt{proc\_bind} no está soportada. El anclaje lo introduce el entorno a
través de una librería nativa. En Finisterrae~III, NumPy carga por debajo
\texttt{libgomp}, y cuando el entorno define \texttt{OMP\_PROC\_BIND} y
\texttt{OMP\_PLACES}, \texttt{libgomp} restringe la máscara de CPU del proceso a
un único núcleo durante el propio \texttt{import} de NumPy, antes de que OMP4Py
cree ningún hilo. Como los intérpretes \emph{free-threaded} heredan la máscara
de CPU del proceso, todos los hilos que crea OMP4Py quedan confinados a ese
núcleo y el cómputo se serializa. El efecto aparece en cualquier modo de
ejecución, porque se produce al importar NumPy y no en una capa concreta de la
biblioteca. Al eliminar esas variables, la máscara abarca todos los núcleos
asignados y EP y FT vuelven a escalar, de forma comparable en Python~3.13t,
3.14t y 3.15t.

Una comprobación controlada lo confirma. Cambiando solo esas variables sobre el
mismo binario y el mismo nodo, EP en modo~3 y clase~S se mantiene en torno a
1.44~s de 1 a 32~hilos cuando están activas y baja a 0.19~s al eliminarlas,
mientras la máscara medida pasa de los 32~núcleos de la asignación a uno solo en
el \texttt{import} de NumPy. El mismo patrón se repite en los cuatro modos de
ejecución. El planificador de Slurm no
define esas variables por su cuenta, de modo que en las campañas preliminares
procedían del entorno de ejecución.

También se realizaron pruebas de aislamiento. \texttt{sys.\_is\_gil\_enabled()}
se mantuvo en \texttt{False} antes y después de importar OMP4Py, Cython y los
módulos compilados. Además, un microbenchmark con hilos Python puros y un
kernel OMP4Py mínimo usaron varias CPUs en Python~3.14t y Python~3.15t. Estas
pruebas apoyan la interpretación de que los problemas restantes pertenecen a
los kernels concretos o a su interacción con el runtime, no a una reactivación
implícita del GIL.


\section{Diagnóstico común con \texttt{perf stat}}
\label{sec:profiling-perf-comun}

\begin{table}[htbp]
\centering
\small
\resizebox{\textwidth}{!}{%
\begin{tabular}{lrrrrrr}
\hline
Bench. & Hilos & NAS (s) & Speedup & CPUs usadas & IPC & Fallos cache (\%) \\
\hline
EP & 1 & 2.87 & 1.00$\times$ & 0.97 & 1.09 & 14.8 \\
EP & 8 & 0.37 & 7.76$\times$ & 2.47 & 1.02 & 44.8 \\
EP & 32 & 0.48 & 5.98$\times$ & 2.55 & 1.09 & 18.9 \\
FT & 1 & 4.90 & 1.00$\times$ & 0.97 & 2.71 & 14.0 \\
FT & 8 & 0.70 & 7.00$\times$ & 3.57 & 2.49 & 30.0 \\
FT & 32 & 0.64 & 7.66$\times$ & 5.72 & 1.88 & 44.2 \\
CG & 1 & 0.78 & 1.00$\times$ & 0.99 & 3.07 & 6.7 \\
CG & 8 & 4.59 & 0.17$\times$ & 2.41 & 1.04 & 3.0 \\
CG & 32 & 6.42 & 0.12$\times$ & 5.98 & 0.39 & 44.5 \\
IS & 1 & 0.03 & 1.00$\times$ & 0.89 & 1.93 & 25.2 \\
IS & 8 & 0.03 & 1.00$\times$ & 1.08 & 1.90 & 23.1 \\
IS & 32 & 0.08 & 0.38$\times$ & 1.23 & 1.83 & 25.8 \\
MG & 1 & 0.46 & 1.00$\times$ & 0.93 & 2.86 & 25.0 \\
MG & 8 & 0.60 & 0.77$\times$ & 3.04 & 0.85 & 39.8 \\
MG & 32 & 0.93 & 0.49$\times$ & 3.95 & 0.65 & 40.7 \\
\hline
\end{tabular}
}
\caption{Perfil homogéneo con \texttt{perf stat} en Finisterrae~III: Python~3.15t, clase~W, modo~3 y recuentos de 1, 8 y 32~hilos para todos los benchmarks.}
\label{tab:profiling-perf-315}
\end{table}


La Tabla~\ref{tab:profiling-perf-315} separa tres comportamientos. EP y FT
reducen el tiempo al añadir hilos. En EP, el tiempo baja de 2.87~s a 0.37~s con
8~hilos, y en FT baja de 4.90~s a 0.70~s. La ocupación de CPU confirma que ambos
kernels ejecutan trabajo paralelo real. FT alcanza 3.57~CPUs con 8~hilos y
5.72~CPUs con 32. En FT, sin embargo, el IPC cae y los fallos de caché suben a
32~hilos, lo que encaja con el óptimo observado alrededor de 16~hilos en la
campaña principal.

CG y MG muestran el patrón contrario. Activan varias CPUs, pero el tiempo
empeora. CG pasa de 0.78~s con un hilo a 4.59~s con 8 y 6.42~s con 32. Al mismo
tiempo, las CPUs usadas suben hasta 5.98 y el IPC cae de 3.07 a 0.39. Por tanto,
el proceso no está serializado, sino que ejecuta paralelismo ineficiente. MG es
menos extremo, pero sigue la misma dirección. Usa 3.04~CPUs con 8~hilos y
3.95~CPUs con 32, mientras el tiempo aumenta y el IPC baja.

IS debe leerse con más cautela porque sus tiempos son de centésimas de segundo.
La señal de rendimiento sigue siendo útil. Con 32~hilos el tiempo sube a
0.08~s, y el trabajo por hilo no compensa la coordinación, los histogramas y las
fases de acumulación. Sus contadores de CPU, en cambio, son menos informativos
porque el arranque del proceso pesa mucho respecto a la ventana del kernel.


\section{Interpretación conjunta}
\label{sec:profiling-interpretacion}

El diagnóstico con \texttt{perf stat} distingue ya tres comportamientos. EP y FT
usan varias CPUs y reducen el tiempo. CG y MG activan varias CPUs, pero el IPC
cae al aumentar los hilos y el tiempo empeora. IS es demasiado corto para
sostener trabajo entre sincronizaciones. La Sección~\ref{sec:profiling-py315}
cruza esta lectura, benchmark a benchmark, con el muestreo y el trazado de
Python~3.15.

La conclusión del diagnóstico es que los malos speedups no tienen una única
causa. En las campañas preliminares, la afinidad heredada del entorno que se
describe en la Sección~\ref{sec:profiling-afinidad} falseaba la lectura de
Python~3.15t. En las finales ese efecto desaparece en EP y FT, y los límites
restantes dependen de la estructura de cada benchmark y de la eficiencia del
runtime en regiones con sincronización o granularidad difícil.


\section{Profiling con el nuevo sistema de Python 3.15}
\label{sec:profiling-py315}

El análisis con \texttt{perf stat} caracteriza el uso de recursos a nivel de
proceso (CPUs, IPC, caché), pero no descompone el tiempo entre trabajo útil y
sincronización dentro del runtime. Python~3.15 incorpora un nuevo sistema de
profiling en la biblioteca estándar, apoyado en la interfaz de depuración externa
de CPython~\cite{PEP768}, con dos componentes. El de muestreo,
\texttt{profiling.sampling}, toma muestras estadísticas de las pilas de llamadas
de todos los hilos del proceso~\cite{PySamplingProfiler}. En modo \emph{wall}, un
hilo bloqueado en una barrera o una reducción aparece en la muestra dentro de su
función de espera, lo que permite atribuir tiempo a la sincronización por nombre
de función. El de trazado, \texttt{profiling.tracing}, es un profiler
determinista que registra cada llamada y devuelve recuentos exactos por función.
Esta sección usa el muestreo para repartir el tiempo entre cómputo y
sincronización, y el trazado para contar cuántas veces se ejecutan las
operaciones del runtime, dos vistas que \texttt{perf stat} solo sugiere de forma
indirecta.

\subsection{Entorno y metodología}
La campaña con el profiler de Python~3.15 se ejecutó en Finisterrae~III con
Python~3.15.0b1t \emph{free-threaded}, clase~S y modo~1. Se escogió modo~1
(híbrido, hilos Python reales sobre el runtime interpretado) porque en ese modo
las primitivas de sincronización de OMP4Py son funciones Python visibles al
profiler. En modo~3 los kernels y sus barreras se compilan a código nativo y el
muestreador de pilas Python no las observa. Cada ejecución de muestreo se lanzó
en modo \emph{wall}, con muestreo de todos los hilos y una frecuencia de
2000~Hz.

El trazado completó los cinco benchmarks con 8~hilos y proporciona los recuentos
estructurales de regiones, hilos creados y barreras. El muestreo aporta una
segunda vista, basada en muestras de pila, para los puntos que completaron de
forma válida en Finisterrae~III. Las campañas de rendimiento de los capítulos
anteriores siguen siendo la referencia para tiempos y speedups. El profiling de
Python~3.15 se usa para identificar patrones estructurales del runtime, como
espera en barreras, reducciones, creación de equipos y reparto de trabajo entre
hilos.

\subsection{Del muestreo a las tablas}
El muestreador imprime una tabla de funciones ordenada por número de muestras.
La obtenida para EP con 8~hilos, recortada en ancho, es la siguiente.

\begin{verbatim}
$ python3.15t -m profiling.sampling run --mode wall -r 2000 \
      examples/EP_Python.py -c S -t 8 -m 1

Profile Stats:
  nsamples  sample%  tottime(ms)  filename:lineno(function)
      1461      7.9        730.5   EP_Python.py:214(...__omp_parallel)
      1351      7.3        675.5   threading.py:363(Condition.wait)
      1331      7.2        665.5   EP_Python.py:210(...__omp_parallel)
      1022      5.5        511.0   EP_Python.py:136(vranlc_ep)
       ...
\end{verbatim}

Las filas dentro del cuerpo de la región (\texttt{\_\_omp\_parallel}) y de
\texttt{vranlc\_ep} son cómputo y suman alrededor del 70\% de las muestras. La
fila de \texttt{threading.Condition.wait} (7.3\%) es la espera en la barrera del
runtime. Para las tablas de este apartado, que necesitan la atribución por hilo,
las ejecuciones se relanzaron con \texttt{--collapsed}, que en lugar de la tabla
agregada vuelca una pila por muestra con su hilo. Un script clasifica entonces la
hoja de cada pila como cómputo, sincronización o inicialización. La fracción de
sincronización de la Tabla~\ref{tab:profiling-py315-sync} sale del cociente entre
las muestras de espera y la suma de espera y cómputo. El número de hilos de la
Tabla~\ref{tab:profiling-py315-churn} es la cantidad de identificadores de hilo
distintos que aparecen en esos volcados.

\subsection{Coste de sincronización}
La Tabla~\ref{tab:profiling-py315-sync} mide la fracción de las muestras de
ejecución paralela que caen en primitivas de espera del runtime. El hilo que
espera aparece en \texttt{threading.Condition.wait}, alcanzada a través de
\texttt{parallelism.py:parallel\_run} y la barrera interna del runtime. En EP,
la espera es muy baja con pocos hilos y sube al 6.9\% con 16~hilos. En CG, FT y
MG la fracción de sincronización ya es mayor con 4~hilos y aumenta al subir a
16~hilos. IS es el caso más extremo. Al tener poco trabajo útil por región, la
espera alcanza casi la mitad del tiempo paralelo con 16~hilos. Este patrón es
coherente con la ley de Amdahl~\cite{Amdahl1967} y con los resultados de
rendimiento. El tipado del modo~3 acelera el cómputo pero no elimina la
coordinación, de modo que la fracción de espera observada aquí señala el límite
efectivo de escalabilidad cuando el kernel es rápido.

\begin{table}[htbp]
\centering
\small
\begin{tabular}{lrrr}
\hline
Benchmark & 1 hilo & 4 hilos & 8 hilos \\
\hline
EP & 0.0\% & 1.0\% & 16.3\% \\
FT & 0.0\% & 9.0\% & 17.4\% \\
CG & 0.1\% & 11.6\% & 19.7\% \\
MG & 0.0\% & 9.7\% & 17.0\% \\
\hline
\end{tabular}
\caption{Coste de sincronización medido con el profiler de muestreo de Python~3.15 (\texttt{profiling.sampling}, modo \emph{wall}, todos los hilos). Porcentaje de las muestras de ejecución paralela que caen en primitivas de espera del runtime OMP4Py (barrera y reducción, \texttt{threading.Condition.wait}), frente al número de hilos. Clase~S, modo~1.}
\label{tab:profiling-py315-sync}
\end{table}


\subsection{Gestión de hilos y balance de carga}
El profiler revela además una diferencia estructural que \texttt{perf stat} no
expone (Tabla~\ref{tab:profiling-py315-churn}). EP ejecuta toda su región
paralela sobre un único equipo persistente de hilos, y con 16~hilos se observan
15~hilos trabajadores distintos. CG, FT, IS y MG, en cambio, crean y destruyen
equipos entre regiones paralelas, de modo que a lo largo de una ejecución
aparecen decenas, cientos o miles de identificadores de hilo, 179 en IS, 512 en
FT, 740 en MG y 6797 en CG. Esa creación y destrucción repetida de equipos es en
sí misma una fuente de sobrecoste y explica parte de la degradación al aumentar
la concurrencia. Sobre el equipo persistente de EP, el reparto de muestras de
cómputo entre hilos trabajadores es prácticamente equilibrado. En los kernels con
rotación de equipos no existe un equipo estable sobre el que medir balance, y el
propio recuento de hilos efímeros pasa a ser el indicador relevante. El origen de
esa rotación está en el diseño del runtime, que se analiza en la
Sección~\ref{sec:profiling-implicaciones}.

\begin{table}[htbp]
\centering
\small
\begin{tabular}{lrr}
\hline
Benchmark & Hilos trabajadores distintos & Sincr./paralelo \\
\hline
EP & 15 (equipo persistente) & 6.9\% \\
FT & 512 & 5.5\% \\
CG & 6797 & 14.2\% \\
IS & 179 & 49.3\% \\
MG & 740 & 10.8\% \\
\hline
\end{tabular}
\caption{Gestión de hilos medida con \texttt{profiling.sampling} en
Finisterrae~III. Hilos trabajadores distintos observados con 16~hilos;
clase~S, modo~1.}
\label{tab:profiling-py315-churn}
\end{table}


\newpage
\subsection{Recuento exacto con el profiler de trazado}
El muestreo reparte el tiempo, pero no cuenta cuántas veces ocurre cada
operación. Ese recuento lo aporta el componente de trazado,
\texttt{profiling.tracing}, un profiler determinista que registra cada llamada.
Su instrumentación infla los tiempos absolutos, así que aquí solo se usan los
recuentos de llamadas, que son exactos. Sobre CG en clase~S, modo~1 y 8~hilos, el
trazado contabiliza 1682~regiones paralelas, 11\,774~creaciones de hilo
(\texttt{threading.Thread}) y 27\,168~entradas en la barrera del runtime
(\texttt{task\_barrier}). El muestreo observa miles de identificadores de hilo,
pero el trazado da el número real, porque captura también los hilos demasiado
breves para aparecer en una muestra. EP, con una sola región paralela, crea su
equipo una única vez y reutiliza 7~hilos durante todo el cómputo. La diferencia
entre 7~hilos y casi 12\,000 pone
cifra al sobrecoste de crear y destruir un equipo por región, y confirma que el
límite de CG está en el runtime, no en el cómputo del kernel.

La Tabla~\ref{tab:profiling-py315-trace} extiende el recuento a los cinco
benchmarks. El número de regiones crece de 1 en EP a 12 en IS, 40 en FT, 127 en
MG y 1682 en CG, y los hilos creados lo acompañan, de 7 a 11\,774. Ese gradiente
coincide con el orden de escalabilidad observado en los capítulos anteriores.

\begin{table}[htbp]
\centering
\small
\begin{tabular}{lrrr}
\hline
Benchmark & Regiones & Hilos creados & Barreras \\
\hline
EP & 1 & 7 & 16 \\
FT & 40 & 280 & 640 \\
CG & 1682 & 11774 & 27168 \\
IS & 12 & 84 & 184 \\
MG & 127 & 889 & 3936 \\
\hline
\end{tabular}
\caption{Recuentos exactos del runtime medidos con \texttt{profiling.tracing}
en Finisterrae~III. Clase~S, modo~1 y 8~hilos.}
\label{tab:profiling-py315-trace}
\end{table}


\subsection{Síntesis por benchmark}
Cruzando las herramientas, cada benchmark tiene un perfil propio. Las medidas
con \texttt{perf stat} corresponden a clase~W y modo~3. Las del muestreo y el
trazado corresponden a clase~S y modo~1 en Finisterrae~III, de modo que se leen
como diagnósticos complementarios y no como una única serie temporal.

EP escala. En clase~W, \texttt{perf} mide entre 3 y 5~CPUs activas y un IPC
estable, sin fallos de caché destacables. En clase~S, el muestreo registra poca
espera a baja concurrencia y el trazado lo explica. Abre una sola región y
reutiliza un equipo de 7~hilos durante todo el cómputo, sin rotación. Su único
límite a recuentos altos es el coste fijo de esa región, no la estructura del
kernel.

FT escala con meseta. \texttt{perf} muestra varias CPUs activas, pero el IPC baja
y los fallos de caché suben a 32~hilos, coherente con una FFT que combina cómputo
y accesos a memoria. El muestreo ve espera por sincronización en las pasadas
paralelas, y el trazado cuenta 40~regiones y 280~hilos creados. El óptimo medido
está en 16~hilos.

CG no escala. \texttt{perf} llega a activar unas 6~CPUs, pero el IPC se desploma
de 3.07 con un hilo a 0.39 con 32 y el tiempo empeora. El muestreo atribuye cerca
del 14\% del tiempo paralelo a espera con 16~hilos, y el trazado descubre la
causa estructural, con 1682~regiones y 11\,774~hilos creados en una sola ejecución de
clase~S. El cuello es la coordinación, no el cómputo.

IS es un caso de trabajo insuficiente. Su kernel es tan corto que \texttt{perf}
resulta ruidoso. El trazado muestra pocas regiones (12) y pocos hilos (84), así
que el problema no es la rotación de equipos, sino que hay tan poco trabajo por
región que el coste de coordinar varios hilos lo supera enseguida.

MG no escala, y combina dos efectos. \texttt{perf} mide 3 o 4~CPUs activas pero
el tiempo sube, con el IPC a la baja y más fallos de caché. El trazado cuenta
127~regiones y 889~hilos creados, una rotación intermedia entre IS y CG, y a
ella se añade la granularidad desigual de los niveles de malla, donde los niveles
gruesos activan un equipo para apenas 1 o 2 iteraciones útiles.

En conjunto, el muestreo y el trazado de Python~3.15 llevan el diagnóstico desde
los síntomas agregados de \texttt{perf stat} hasta una atribución por función,
hilo y recuento, que es la base de las mejoras de runtime de la
Sección~\ref{sec:profiling-implicaciones}.


\section{Implicaciones para OMP4Py}
\label{sec:profiling-implicaciones}

El diagnóstico no invalida el port. EP y FT muestran que las regiones paralelas
adaptadas con OMP4Py producen speedups altos cuando el kernel es adecuado, y el
modo~3 confirma que el tipado Cython elimina gran parte del coste de objetos
Python en los bucles críticos. Más allá del comportamiento de cada benchmark, el
nuevo sistema de profiling de Python~3.15 señala dos puntos del propio runtime
que limitan la escalabilidad con independencia del kernel, y que son los
candidatos más claros a mejora en la biblioteca. El análisis se hizo sobre el
runtime interpretado. Como el runtime nativo conserva la misma estructura general
de equipos, barreras y reducciones, es razonable esperar que parte de estos
límites se mantenga en los modos compilados. La magnitud exacta, sin embargo,
requiere instrumentación interna del runtime nativo.

\subsection{Reutilización del equipo de hilos}
El primer límite es la gestión del equipo de hilos. En su versión actual, el
runtime crea un equipo nuevo en cada región paralela. La función
\texttt{parallel\_run} lanza un \texttt{threading.Thread} por trabajador y lo
descarta al cerrar la región, sin un pool persistente que se reutilice entre
regiones. El coste de crear y destruir hilos es fijo por región, de modo que un
kernel con una sola región (EP) lo paga una vez, mientras que un kernel con
muchas regiones cortas (CG, o el ciclo en V de MG) lo paga miles de veces. Es la
diferencia entre el equipo estable de EP y los miles de identificadores de hilo
que aparecen en CG (Tabla~\ref{tab:profiling-py315-churn}). De hecho, los
benchmarks que
escalan se reescribieron para envolver todo el cómputo en una única región
\texttt{parallel} y repartir el trabajo interno con \texttt{for}, justamente para
evitar esa rotación. Que esa reescritura sea necesaria indica que la mejora
corresponde al runtime. Un equipo persistente, creado una vez y aparcado en una
barrera entre regiones, eliminaría el sobrecoste sin obligar al usuario a
reestructurar su código.

\subsection{Coste de la barrera y la reducción}
El segundo límite es el coste de la sincronización. Cada región termina en una
barrera (\texttt{task\_barrier}) construida sobre \texttt{threading.Event}, que a
su vez bloquea en \texttt{threading.Condition.wait}. Esas primitivas de la
biblioteca estándar son relativamente pesadas, y su coste se nota cuando el
trabajo entre barreras es pequeño. La fracción de tiempo en espera alcanza
valores relevantes con 16~hilos, especialmente en CG, IS y MG
(Tabla~\ref{tab:profiling-py315-sync}).
Como cada reducción escalar conlleva además una barrera, un kernel con varias
reducciones por iteración, como las cuatro de CG, acumula ese coste. Una barrera
más ligera, por ejemplo una barrera centralizada con contador o una
implementación con operaciones atómicas en el runtime nativo, y una reducción que
no exija una barrera completa, reducirían esa fracción de espera.

\subsection{Calidad del código compilado}
El tercer nivel es la calidad del código que genera el modo compilado. Si en el
bucle crítico queda una llamada a una función Python externa, Cython no puede
compilarla y el tiempo absoluto vuelve a estar dominado por el intérprete, como
ocurría en EP con \texttt{vranlc} antes de incluir una versión compilable por Cython
(Sección~\ref{sec:decisiones-implementacion}). Mantener todo el cuerpo del kernel
compilable es, por tanto, tan importante como el tipado.

Estas observaciones provienen de profiling externo sobre el runtime interpretado
en Finisterrae~III. El siguiente paso sigue siendo instrumentar OMP4Py
internamente, por región y para los modos compilados, midiendo cuántos hilos
entran en cada región y qué fracción de tiempo pasan en trabajo útil frente a
creación de equipos, barreras y reducciones. Eso convertiría estas implicaciones
en una guía cuantitativa para el desarrollo de la biblioteca.
